\begin{frame}{1.2 마무리}
  \begin{block}{핵심 질문}
    어떤 문제를 만나든, 첫 질문은 항상 같다:
    \kq{손에 쥔 데이터에 무엇이 적혀 있는가 --- 그리고 그것을
    ``라벨''으로 볼 것인가, ``보상''으로 볼 것인가, 아니면
    ``구조''로만 볼 것인가.}
  \end{block}
  \vskip0.6em
  \begin{itemize}
    \item $x, y, h$ 로 \textbf{먼저 정의}하고, 두 질문으로 세 갈래를
      가르고, 지도 안에서는 ``연속/범주''로 회귀/분류를 가른다.
    \item 회귀 vs 분류 = \textbf{같은 모델, 다른 출력}
      (항등 링크 vs 시그모이드).
    \item 모든 알고리즘의 뼈대 = \textbf{$J$ 를 정의하고 최소화}
      ($h$, $J$, 최소화 방법의 세 질문).
  \end{itemize}
\end{frame}
